% =====================================================================
%  CLASE 07 -- Drenaje subterráneo: régimen permanente
% =====================================================================
\clase{07}{Drenaje subterráneo: régimen permanente}

\section{Motivación}
Suelos agrícolas anegados, excavaciones inundadas por la napa o canchas
deportivas saturadas: ¿cómo podemos solucionar este problema?

\section{Drenaje subterráneo}
Los drenes son utilizados con distintos fines:
\begin{itemize}
  \item Disminuir los niveles de agua para evitar la saturación de las raíces
        de una plantación.
  \item Evitar la saturación superficial del suelo de complejos deportivos.
  \item Captación de aguas subterráneas con fines de almacenamiento de agua
        potable y/o superficial.
\end{itemize}
Tipos de obras:
\begin{itemize}
  \item Pozos.
  \item Drenes:
  \begin{itemize}
    \item Drenes verticales o punteras.
    \item Drenes horizontales: naturales, interceptores y paralelos.
  \end{itemize}
\end{itemize}

\section{Drenes horizontales}
\begin{itemize}
  \item Conocidos como drenes o galerías de infiltración.
  \item No son tan populares y usados como los pozos.
  \item Usados en napas superficiales de poco espesor.
  \item Permiten eliminar el consumo energético de bombas.
\end{itemize}

\subsection{Drenes interceptores}
\begin{itemize}
  \item Evacúan sólo una parte del flujo.
  \item $h_0$ es la altura en el canal de evacuación.
  \item Dimensionar $h_0$ tal que el problema se solucione.
  \item $h_0$ puede tomarse como la altura normal en el canal (dren interceptor).
\end{itemize}

\begin{figure}[H]
\centering
\begin{tikzpicture}[font=\small]
  \fill[pattern=crosshatch,pattern color=gray!60] (0,-0.25) rectangle (9,0);
  \draw (0,0) -- (9,0);
  % canal
  \draw[thick] (3,2.4) -- (3,0) -- (4.6,0) -- (4.6,2.4);
  \fill[agua!50] (3,0) rectangle (4.6,0.9);
  \pic at (4.2,0.92) {nivelagua};
  % napa aguas arriba
  \draw[azulusm,thick] (4.6,0.9) .. controls (5.3,1.0) and (5.6,2.0) .. (6.5,2.1) -- (9,2.15);
  \draw[azulusm,thick] (0,0.8) -- (3,0.9);
  \draw[cota] (3.3,0) -- (3.3,0.9) node[midway,right]{$h_0$};
  \draw[cota] (5.6,0) -- (5.6,1.75) node[midway,right]{$h$};
  \draw[cota] (8.2,0) -- (8.2,2.15) node[midway,right]{$H$};
  \draw[cota] (3,2.6) -- (4.6,2.6) node[midway,above]{$b$};
  \draw[->] (4.6,-0.5) -- (5.6,-0.5) node[midway,below]{$x$};
  \draw[->] (4.6,-1.1) -- (8.2,-1.1) node[midway,below]{$x_1$};
  \draw[flecha,azulusm] (7.6,1.2) -- (6.8,1.2) node[pos=0,above]{$q$};
  \draw[flecha,azulusm] (2.2,0.4) -- (1.2,0.4) node[pos=0,above]{$q'$};
\end{tikzpicture}
\caption{Dren interceptor (corte transversal): el canal capta el caudal $q$
y deja pasar $q'$.}
\end{figure}

Aplicando Dupuit al flujo hacia el dren:
\begin{formula}
$\displaystyle h^2-h_0^2=\frac{2\cdot q}{K}\cdot(x-x_0)$\\[4pt]
Si para $x=x_1$, $h=H$ y $x_0=0$:\qquad $\displaystyle H^2-h_0^2=\frac{2\cdot q}{K}\cdot x_1$
\end{formula}
Lo que se lleva el canal, con $L$ el largo del dren:
\begin{formula}
$Q_z=(q-q')\cdot L$
\end{formula}

\subsection{Drenes paralelos}
\paragraph{Clasificación según ubicación en el acuífero}
\begin{itemize}
  \item \textbf{Zanjas abiertas:}
  \begin{itemize}
    \item Líneas de flujo casi horizontales, excepto cerca de la captación.
    \item Equipotenciales planas cercanas a la vertical.
  \end{itemize}
  \item \textbf{Tuberías:}
  \begin{itemize}
    \item Líneas de flujo radiales dirigidas a la captación.
    \item Equipotenciales son superficies semicilíndricas concéntricas con la captación.
  \end{itemize}
\end{itemize}

\begin{figure}[H]
\centering
\begin{subfigure}{0.48\textwidth}\centering
\begin{tikzpicture}[font=\small,scale=0.8]
  \fill[aguaclara] (0,0) rectangle (6,3);
  \draw[thick] (0,3.3) -- (6,3.3);
  \draw[thick] (2.9,0) -- (2.9,3.3) (3.1,0) -- (3.1,3.3);
  \draw[azulusm,thick] (0,2.9) .. controls (2,2.8) and (2.7,2.5) .. (2.9,2.2);
  \draw[azulusm,thick] (6,2.9) .. controls (4,2.8) and (3.3,2.5) .. (3.1,2.2);
  \foreach \y in {0.5,1.2,1.9} {\draw[->,azulusm] (0.3,\y) -- (2.8,\y); \draw[->,azulusm] (5.7,\y) -- (3.2,\y);}
  \foreach \x in {0.8,1.6,2.3,3.7,4.4,5.2} \draw[dashed] (\x,0) -- (\x,2.8);
  \fill[pattern=north east lines] (0,-0.2) rectangle (6,0);
\end{tikzpicture}
\caption{Zanjas abiertas.}
\end{subfigure}\hfill
\begin{subfigure}{0.48\textwidth}\centering
\begin{tikzpicture}[font=\small,scale=0.8]
  \fill[aguaclara] (0,0) rectangle (6,3);
  \draw[thick] (0,3.3) -- (6,3.3);
  \draw[azulusm,thick] (0,2.9) .. controls (2,2.8) and (2.6,2.5) .. (3,2.3) .. controls (3.4,2.5) and (4,2.8) .. (6,2.9);
  \fill[azulusm] (3,2.3) circle (0.18);
  \foreach \a in {0,20,...,180} \draw[->,azulusm] (3,2.3) ++(-\a:2.0) -- ++(180-\a:1.7);
  \foreach \r in {0.8,1.4,2.0} \draw[dashed] (3,2.3) ++(0:\r) arc (0:-180:\r);
  \fill[pattern=north east lines] (0,-0.2) rectangle (6,0);
\end{tikzpicture}
\caption{Tuberías.}
\end{subfigure}
\caption{Redes de flujo hacia drenes paralelos.}
\end{figure}

\paragraph{Régimen de escurrimiento}
\begin{itemize}
  \item Condición de equilibrio (régimen permanente).
  \item Condición de desequilibrio (régimen impermanente).
\end{itemize}

% Macro: drenes paralelos abiertos (#1 = separación al fondo impermeable d)
\newcommand{\drenesabiertos}[1]{%
\begin{tikzpicture}[font=\small,scale=0.9]
  \def\d{#1}
  \fill[pattern=crosshatch,pattern color=gray!60] (-0.5,-\d-0.25) rectangle (8.5,-\d);
  \draw (-0.5,-\d) -- (8.5,-\d);
  % zanjas
  \draw[thick] (-0.5,3) -- (0.6,3) -- (0.6,0) -- (1.0,0) -- (1.0,3) -- (7.0,3) -- (7.0,0) -- (7.4,0) -- (7.4,3) -- (8.5,3);
  \fill[agua!60] (0.6,0) rectangle (1.0,1.0); \fill[agua!60] (7.0,0) rectangle (7.4,1.0);
  % volumen de control
  \fill[rojo,opacity=0.15,draw=rojo] (1.0,0) rectangle (3.3,3);
  \node[rojo] at (1.6,2.6) {\large$V_c$};
  % napa
  \draw[azulusm,thick] (1.0,1.0) .. controls (1.3,2.0) and (2.5,2.3) .. (4.0,2.3) .. controls (5.5,2.3) and (6.7,2.0) .. (7.0,1.0);
  \pic at (0.8,1.02) {nivelagua}; \pic at (7.2,1.02) {nivelagua};
  % lluvia
  \foreach \x in {-0.3,0.1,1.6,2.1,2.6,3.1,3.6,4.1,4.6,5.1,5.6,6.1,7.8,8.2}
    \draw[->,azulusm,thick] (\x,4.0) -- (\x,3.1);
  \node at (0.0,4.3) {$f$}; \node at (3.8,4.3) {$f$}; \node at (8.0,4.3) {$f$};
  \draw[cota] (0.2,0) -- (0.2,1.0) node[midway,left]{$h_0$};
  \draw[cota] (7.8,0) -- (7.8,1.0) node[midway,right]{$h_0$};
  \draw[cota] (4.0,0) -- (4.0,2.3) node[midway,right]{$H$};
  \draw[cota] (1.3,0) -- (1.3,1.45) node[midway,right]{$h$};
  \draw[->] (1.0,0.15) -- (2.0,0.15) node[midway,above]{$x$};
  \draw[flecha,azulusm] (2.4,1.1) -- (1.7,1.1) node[pos=0,right]{$q$};
  \draw[flecha,azulusm] (5.8,1.1) -- (6.5,1.1) node[pos=0,left]{$q$};
  \draw[cota] (0.8,-\d-0.6) -- (7.2,-\d-0.6) node[midway,below]{$D$};
  \ifdim\d pt>0pt \draw[cota] (7.6,-\d) -- (7.6,0) node[midway,right]{$d$};\fi
\end{tikzpicture}}

\section{Drenes abiertos: drenes paralelos en condición de equilibrio}

\subsection{Drenes que abarcan toda la napa: ecuación de Donnan}
\begin{itemize}
  \item Acuífero homogéneo e isotrópico.
  \item Drenes abarcan toda la napa (llegan al fondo impermeable).
  \item Espaciados equidistantemente.
  \item Alimentación constante por unidad de superficie ($f$).
  \item Se desprecia la curvatura del flujo cerca de la zanja.
\end{itemize}

\begin{figure}[H]
\centering
\drenesabiertos{0}
\caption{Drenes paralelos abiertos que abarcan toda la napa.}
\end{figure}

\begin{nota}[title={Deducción}]
Balance en el volumen de control $V_c$ entre el dren ($x=0$) y la sección
$x$, y Dupuit: el caudal que sale por $x=0$ es $q=f\left(\frac D2-x\right)$ y
\[
q=K\,h\frac{\dd h}{\dd x}\ \Rightarrow\ f\left(\frac D2-x\right)\dd x=K\,h\,\dd h
\ \Rightarrow\ \int_0^{D/2}f\left(\frac D2-x\right)\dd x=\int_{h_0}^{H}K\,h\,\dd h
\]
\end{nota}
\begin{formula}
$\displaystyle H^2-h_0^2=\frac{f\cdot D^2}{4\cdot K}$\qquad\textbf{Ecuación de Donnan}
\end{formula}

\subsection{Drenes que no abarcan toda la napa: ecuación de Hooghoudt}
\begin{itemize}
  \item Acuífero homogéneo e isotrópico.
  \item Drenes \textbf{no} abarcan toda la napa.
  \item Espaciados equidistantemente.
  \item Alimentación constante por unidad de superficie.
  \item La distancia al fondo impermeable no debe ser muy grande ($d<0.6$ m).
  \item Se desprecia la curvatura del flujo cerca de la zanja.
\end{itemize}

\begin{figure}[H]
\centering
\drenesabiertos{0.8}
\caption{Drenes paralelos abiertos a una distancia $d$ del estrato impermeable.}
\end{figure}

\begin{formula}
$\displaystyle(H-h_0)(H+h_0+2d)=\frac{f\cdot D^2}{4\cdot K}$\qquad\textbf{Ecuación de Hooghoudt}
\end{formula}

\paragraph{Corrección de Hooghoudt.} Para $d>0.6$ m (flujo radial hacia el
dren) se utiliza una profundidad equivalente $d'$, que se obtiene de un
ábaco en función de la profundidad del estrato $d$ y del espaciamiento $D$
($d'<d$, y $d'$ tiende a un valor límite al crecer $d$):
\begin{formula}
$\displaystyle(H-h_0)(H+h_0+2d')=\frac{f\cdot D^2}{4\cdot K}$
\end{formula}

\section{Drenes cerrados: drenes paralelos en condición de equilibrio}

% Macro: drenes cerrados (tuberías)
\newcommand{\drenescerrados}[1]{% #1: 1 = fondo a distancia d, 0 = infinito
\begin{tikzpicture}[font=\small,scale=0.9]
  \foreach \x in {0,0.4,...,8} \draw[->,azulusm,thick] (\x,3.4) -- (\x,2.6);
  \node at (0.6,3.7) {$f$}; \node at (4,3.7) {$f$}; \node at (7.4,3.7) {$f$};
  \draw[thick] (-0.2,2.5) -- (8.2,2.5);
  \draw[dashed] (-0.5,1.0) -- (8.5,1.0);
  \foreach \x in {0.4,4.0,7.6} \draw[thick,fill=white] (\x,1.0) circle (0.35);
  \draw[azulusm,thick] (0.75,1.1) .. controls (1.2,2.1) and (3.2,2.1) .. (3.65,1.1);
  \draw[azulusm,thick] (4.35,1.1) .. controls (4.8,2.1) and (6.8,2.1) .. (7.25,1.1);
  \draw[cota] (2.2,1.0) -- (2.2,1.85) node[midway,right]{$H$};
  \draw[cota] (0.05,1.6) -- (0.75,1.6) node[midway,above]{$2r_0$};
  \draw[cota] (0.4,0.3) -- (4.0,0.3) node[midway,below]{$D$};
  \draw[cota] (4.0,0.3) -- (7.6,0.3) node[midway,below]{$D$};
  \ifnum#1=1
    \fill[pattern=crosshatch,pattern color=gray!60] (-0.5,-0.9) rectangle (8.5,-0.65);
    \draw (-0.5,-0.65) -- (8.5,-0.65);
    \draw[cota] (8.2,-0.65) -- (8.2,1.0) node[midway,right]{$d$};
    \fill[rojo,opacity=0.15,draw=rojo] (4.0,-0.65) rectangle (5.8,2.5);
    \node[rojo] at (4.9,2.2) {$V_c$};
  \else
    \fill[rojo,opacity=0.15,draw=rojo] (2.2,2.5) -- (2.2,1.0) arc (180:360:1.8) -- (5.8,2.5) -- cycle;
    \node[rojo] at (4.0,2.2) {$V_c$};
    \foreach \a in {200,230,...,340} \draw[->,azulusm] (4,1.0) ++(\a:1.7) -- ++(\a+180:1.2);
    \node at (7.5,-0.7) {$d\approx\infty$};
  \fi
\end{tikzpicture}}

\subsection{Corrección de Dagan (1964)}
\begin{itemize}
  \item Para distancias ($d$) al estrato impermeable pequeñas.
  \item Incorpora la curvatura cerca del dren.
  \item Se aplica un factor de corrección a la fórmula de Donnan.
  \item También corrige la fórmula de Donnan ($d=h_0$).
\end{itemize}

\begin{figure}[H]
\centering
\drenescerrados{1}
\caption{Drenes paralelos cerrados (tuberías) a una distancia $d$ del estrato impermeable.}
\end{figure}

\begin{formula}
$\displaystyle(H+d)^2-d^2=\frac{f\cdot D^2}{4\cdot K}\left(1+\alpha\frac{4\cdot d}{D}\right)$\\[8pt]
$\displaystyle\alpha=-\frac1\pi\ln\left[2\left(\cosh\left(\frac{\pi\cdot r_0}{d}\right)-1\right)\right]$
\end{formula}

\begin{table}[H]
\centering
\begin{tabular}{cc}
\toprule
$r_0/d$ & $\alpha$\\
\midrule
0.01 & 2.16\\
0.02 & 1.75\\
0.04 & 1.31\\
0.06 & 1.05\\
0.08 & 0.86\\
0.10 & 0.75\\
0.20 & 0.29\\
0.30 & 0.01\\
\bottomrule
\end{tabular}
\caption{Factor $\alpha$ de Dagan. Para $r_0/d\geq0.3\Rightarrow\alpha\to0$.}
\end{table}

\subsection{Drenes superficiales (acuífero de espesor infinito)}
\begin{itemize}
  \item Acuífero homogéneo e isotrópico.
  \item Acuífero de espesor muy grande o indefinido (infinito).
  \item Espaciados equidistantemente.
  \item Alimentación constante por unidad de superficie.
  \item Escurrimiento semi-cilíndrico.
\end{itemize}

\begin{figure}[H]
\centering
\drenescerrados{0}
\caption{Drenes superficiales en acuífero de espesor infinito: escurrimiento semicilíndrico.}
\end{figure}

\begin{formula}
$\displaystyle H=\frac{D\cdot f}{K\cdot\pi}\ln\left(\frac{D}{5.44\cdot r_0}\right)$
\end{formula}

\subsection{Acuíferos de espesor finito: fórmula de Kirkham (1958)}
Incorpora la curvatura cerca del dren:
\begin{formula}
$\displaystyle H=\frac{D\cdot f}{K}\cdot F\left(\frac{2r_0}{D},\frac dD\right)$\\[8pt]
$\displaystyle F(x,y)=\frac1\pi\left[\ln\left(\frac{2}{\pi x}\right)+\sum_{m=1}^\infty\frac1m\big(\cos(m\pi x)-\cos(m\pi)\big)\big(\coth(2m\pi y)-1\big)\right]$
\end{formula}

\begin{table}[H]
\centering
\begin{tabular}{c|ccccc}
\toprule
& \multicolumn{5}{c}{$2r_0/D$}\\
$d/D$ & 0.0025 & 0.005 & 0.01 & 0.02 & 0.04\\
\midrule
0.01 & 12.79 & 12.57 & 12.33 & 12.03 & 11.52\\
0.02 & 6.761 & 6.541 & 6.318 & 6.077 & 5.771\\
0.04 & 3.864 & 3.643 & 3.421 & 3.195 & 2.954\\
0.08 & 2.522 & 2.301 & 2.080 & 1.858 & 1.633\\
0.16 & 1.961 & 1.741 & 1.520 & 1.299 & 1.077\\
0.32 & 1.787 & 1.566 & 1.345 & 1.125 & 0.904\\
0.64 & 1.764 & 1.543 & 1.323 & 1.102 & 0.811\\
1.00 & 1.763 & 1.543 & 1.322 & 1.101 & 0.811\\
$\infty$ & 1.763 & 1.543 & 1.322 & 1.101 & 0.811\\
\bottomrule
\end{tabular}
\caption{Valores de $F(2r_0/D,\,d/D)$ de Kirkham.}
\end{table}

\section{Régimen permanente: consideraciones de diseño}
La infiltración siempre dependerá de:
\begin{itemize}
  \item Nivel de saturación del suelo.
  \item Intensidad de precipitación.
  \item Propiedades del suelo.
  \item Recordar métodos de infiltración de hidrología para valores precisos.
\end{itemize}
\begin{table}[H]
\centering
\begin{tabular}{lc}
\toprule
\textbf{Tipo de suelo} & $f$ [mm/h]\\
\midrule
Arena & $>30$\\
Arena limosa & 20 -- 30\\
Limo & 10 -- 20\\
Arcilla limosa & 5 -- 10\\
Arcilla & 1 -- 5\\
\bottomrule
\end{tabular}
\caption{Ejemplos de valores de infiltración.}
\end{table}
El diseño en régimen de equilibrio o permanente se utiliza cuando no es
permisible tener un nivel de saturación prolongado en el tiempo. Ejemplo:
plantas con raíces que no permiten estar saturadas por un tiempo prolongado,
campos deportivos, etc.
